% @file paper63_galois_representations_en.tex
% @brief Galois Representations and the Langlands Programme
% @author Michael Fuhrmann
% @date 2026-03-12
% @project Specialist Mathematics Project, Build 143

\documentclass[a4paper,12pt]{amsart}

\usepackage{amsmath, amssymb, amsthm}
\usepackage{hyperref}
\usepackage{geometry}
\usepackage{enumitem}
\geometry{margin=2.5cm}

% Theorem environments
\newtheorem{theorem}{Theorem}[section]
\newtheorem{lemma}[theorem]{Lemma}
\newtheorem{proposition}[theorem]{Proposition}
\newtheorem{corollary}[theorem]{Corollary}
\newtheorem{conjecture}[theorem]{Conjecture}
\theoremstyle{definition}
\newtheorem{definition}[theorem]{Definition}
\newtheorem{example}[theorem]{Example}
\theoremstyle{remark}
\newtheorem{remark}[theorem]{Remark}

% Mathematical shortcuts
\newcommand{\Q}{\mathbb{Q}}
\newcommand{\Z}{\mathbb{Z}}
\newcommand{\Zl}{\mathbb{Z}_\ell}
\newcommand{\Ql}{\mathbb{Q}_\ell}
\newcommand{\Fp}{\mathbb{F}_p}
\newcommand{\Fq}{\mathbb{F}_q}
\newcommand{\Qbar}{\overline{\mathbb{Q}}}
\newcommand{\GQ}{\mathrm{Gal}(\overline{\mathbb{Q}}/\mathbb{Q})}
\newcommand{\GL}{\mathrm{GL}}
\newcommand{\SL}{\mathrm{SL}}
\newcommand{\Frob}{\mathrm{Frob}}
\newcommand{\tr}{\mathrm{tr}}
\newcommand{\rho}{\varrho}
\newcommand{\BdR}{B_{\mathrm{dR}}}
\newcommand{\Bcrys}{B_{\mathrm{crys}}}
\newcommand{\Tl}{T_\ell}
\newcommand{\Vl}{V_\ell}

\title{Galois Representations and the Langlands Programme}

\author{Michael Fuhrmann}
\address{Specialist Mathematics Project, Build 143}
\email{specialist-maths@project.local}
\date{2026-03-12}

\begin{abstract}
Galois representations are continuous group homomorphisms
$\varrho\colon \mathrm{Gal}(\overline{\mathbb{Q}}/\mathbb{Q}) \to \GL_n(R)$,
where $R$ is typically a ring of $\ell$-adic integers or a finite field.
They form the central bridge between algebraic number theory and analytic objects
such as modular forms and automorphic representations, and are the fundamental
language of the Langlands programme.

This paper surveys the theory of Galois representations from the absolute Galois
group through $\ell$-adic and $p$-adic representations to their deep connections
with elliptic curves, modular forms, and the Langlands correspondence.
We discuss the Modularity Theorem of Wiles--Taylor--Wiles (1995), which is proved
and implies Fermat's Last Theorem; Serre's Modularity Conjecture, proved by
Khare--Wintenberger (2009); Fontaine's $p$-adic Hodge theory; and Lafforgue's proof
of the Langlands correspondence over function fields (2002). The global Langlands
correspondence for $\mathbb{Q}$ remains a central open conjecture of mathematics.
\end{abstract}

\maketitle

\tableofcontents

%=========================================================
\section{Introduction}
%=========================================================

The theory of Galois representations stands at the crossroads of two great streams
of modern mathematics: \emph{algebraic} (Galois theory, algebraic number fields,
elliptic curves, motives) and \emph{analytic} (automorphic forms, $L$-functions,
harmonic analysis on Lie groups). A Galois representation encodes the action of
the absolute Galois group $\GQ$ on a finite-dimensional vector space, translating
the intrinsically infinite, topological group into concrete linear algebra.

The driving philosophy of the \textbf{Langlands programme}, formulated by Robert
Langlands in 1967 and 1970, is the principle of \emph{functoriality}: that
$n$-dimensional Galois representations should correspond bijectively to certain
automorphic representations of $\GL_n(\mathbb{A}_\Q)$, the adelic general linear
group. This correspondence, if fully established, would unify vast swaths of
mathematics and yield unprecedented control over $L$-functions.

\subsection*{Historical context}

The roots of Galois representations lie in classical reciprocity laws. Gauss's
quadratic reciprocity (1796) can be rephrased as a statement about the splitting
of primes in quadratic fields, controlled by a 1-dimensional Galois representation
(a Dirichlet character). Class field theory, developed by Hilbert, Takagi, and
Artin through the first third of the twentieth century, classifies all abelian
extensions of number fields and establishes abelian Langlands in the guise of
Artin's reciprocity law.

The non-abelian programme began with Eichler--Shimura (1950s), who showed that the
$L$-function of a weight-2 newform equals the $L$-function of an elliptic curve.
Serre and Deligne systematically developed $\ell$-adic representations in the 1960s
and 1970s. The spectacular culmination came with Wiles's proof of Fermat's Last
Theorem (1995), using the Modularity Theorem for semi-stable elliptic curves, which
was later extended to all elliptic curves over $\Q$ by Breuil--Conrad--Diamond--Taylor.

\subsection*{Structure of this paper}

Sections~\ref{sec:galois-group} and~\ref{sec:ell-adic} set up the absolute Galois
group and $\ell$-adic representations. Section~\ref{sec:elliptic-curves} constructs
representations from elliptic curves via the Tate module. Sections~\ref{sec:modularity}
and~\ref{sec:modular-forms} treat the Modularity Theorem and modular Galois
representations. Sections~\ref{sec:serre} and~\ref{sec:fontaine} cover Serre's
Modularity Theorem and Fontaine's $p$-adic Hodge theory. Finally,
Sections~\ref{sec:langlands} and~\ref{sec:recent} survey the Langlands correspondence
and recent results.

%=========================================================
\section{The Absolute Galois Group}
\label{sec:galois-group}
%=========================================================

\begin{definition}
The \textbf{absolute Galois group} of $\Q$ is
\[
  G_\Q := \GQ = \mathrm{Aut}(\Qbar/\Q),
\]
the group of all field automorphisms of a fixed algebraic closure $\Qbar$ that fix
$\Q$ pointwise.
\end{definition}

\subsection{Profinite topology}

Although $G_\Q$ is enormous (uncountable), it carries a natural topology making it
a profinite group:
\[
  G_\Q \;\cong\; \varprojlim_{[K:\Q] < \infty,\; K/\Q \text{ Galois}} \mathrm{Gal}(K/\Q).
\]
Here the inverse limit is taken over the directed system of finite Galois extensions
$K/\Q$, ordered by inclusion. The topology is the \emph{Krull topology}: a basis of
open subgroups consists of the normal subgroups $\mathrm{Gal}(\Qbar/K)$ for finite
Galois $K/\Q$.

\begin{proposition}
$G_\Q$ is compact, totally disconnected, and Hausdorff (i.e., profinite).
\end{proposition}

\begin{proof}
By Tychonoff's theorem, the product $\prod_K \mathrm{Gal}(K/\Q)$ over all finite
Galois extensions is compact. The inverse limit is a closed subgroup, hence compact.
Total disconnectedness follows because every open subgroup has finite index and the
open subgroups form a neighbourhood base.
\end{proof}

\subsection{Frobenius elements}

For a prime $p$ unramified in a Galois extension $K/\Q$, there is a well-defined
conjugacy class $\Frob_p \in \mathrm{Gal}(K/\Q)$ characterised by
\[
  \Frob_p(\alpha) \equiv \alpha^p \pmod{\mathfrak{P}}
\]
for all algebraic integers $\alpha \in \mathcal{O}_K$ and any prime $\mathfrak{P}$
of $K$ above $p$.

\begin{definition}
Let $\ell$ be a prime. The \textbf{geometric Frobenius} at $p \neq \ell$ is the
element $\Frob_p \in G_\Q$ (or rather the conjugacy class it defines) acting on the
$\ell$-power roots of unity by $\zeta_{\ell^n} \mapsto \zeta_{\ell^n}^p$.
\end{definition}

By the Chebotarev density theorem, the Frobenius elements $\{\Frob_p : p \text{ unramified}\}$
are \emph{dense} in $G_\Q$. Consequently, a continuous representation is completely
determined by the values on Frobenius elements.

\subsection{Ramification}

For each prime $p$, fix an embedding $\Qbar \hookrightarrow \overline{\Q}_p$.
This gives a decomposition group $D_p \subset G_\Q$ and an inertia group $I_p \trianglelefteq D_p$
sitting in the exact sequence
\[
  1 \to I_p \to D_p \to G_{\Fp} \cong \hat{\Z} \to 1.
\]
The arithmetic Frobenius generator of $G_{\Fp}$ lifts to $\Frob_p \in D_p/I_p$.

%=========================================================
\section{\texorpdfstring{$\ell$}{l}-adic Representations}
\label{sec:ell-adic}
%=========================================================

Fix a prime $\ell$. The ring $\Zl$ of $\ell$-adic integers is the completion of $\Z$
at $\ell$, and $\Ql = \Zl[\ell^{-1}]$ is the field of $\ell$-adic numbers.

\begin{definition}
An \textbf{$\ell$-adic representation} of $G_\Q$ of dimension $n$ is a continuous
homomorphism
\[
  \varrho\colon G_\Q \to \GL_n(\Zl),
\]
where $\GL_n(\Zl)$ carries the $\ell$-adic topology. Often we consider the
associated rational representation $\varrho\colon G_\Q \to \GL_n(\Ql)$.
\end{definition}

\subsection{Continuity and finiteness}

Because $G_\Q$ is profinite and $\GL_n(\Zl)$ is an $\ell$-adic Lie group, continuity
means that the kernel contains an open subgroup, equivalently that $\varrho$ factors
through a finite-index open subgroup of $G_\Q$.

\subsection{The cyclotomic character}

The fundamental 1-dimensional $\ell$-adic representation is the \textbf{cyclotomic
character}
\[
  \chi_\ell\colon G_\Q \to \Zl^\times,
\]
defined by $\sigma(\zeta) = \zeta^{\chi_\ell(\sigma)}$ for all $\ell^n$-th roots of
unity $\zeta$. Concretely, via the canonical isomorphism
$\hat{\Z} \cong \varprojlim \Z/\ell^n\Z$, one has
\[
  \chi_\ell(\Frob_p) = p \quad \text{for all primes } p \neq \ell.
\]

\begin{example}[Dirichlet characters]
A Dirichlet character $\chi\colon (\Z/N\Z)^\times \to \mathbb{C}^\times$ of conductor
$N$ gives a 1-dimensional Galois representation via class field theory:
\[
  \varrho_\chi\colon G_\Q \twoheadrightarrow \mathrm{Gal}(\Q(\zeta_N)/\Q) \cong (\Z/N\Z)^\times
  \xrightarrow{\chi} \mathbb{C}^\times.
\]
The Frobenius value is $\varrho_\chi(\Frob_p) = \chi(p)$ for $p \nmid N$.
\end{example}

\subsection{Tate twists}

For $n \in \Z$, the $n$-fold \textbf{Tate twist} of a representation $V$ is
$V(n) := V \otimes \chi_\ell^{\otimes n}$. The category of $\ell$-adic representations
is closed under direct sums, tensor products, duals, and Tate twists.

%=========================================================
\section{Representations from Elliptic Curves}
\label{sec:elliptic-curves}
%=========================================================

Elliptic curves provide the most important source of 2-dimensional Galois representations.

\begin{definition}
Let $E/\Q$ be an elliptic curve and $\ell$ a prime. The \textbf{$\ell$-adic Tate
module} is
\[
  \Tl(E) := \varprojlim_n E[\ell^n],
\]
where $E[\ell^n] \cong (\Z/\ell^n\Z)^2$ are the $\ell^n$-torsion points (over $\Qbar$)
and the maps are multiplication by $\ell$. As an abelian group,
$\Tl(E) \cong \Zl^2$.
\end{definition}

Since $G_\Q$ permutes the torsion points, it acts on $\Tl(E)$, giving

\begin{definition}
The \textbf{$\ell$-adic representation of $E$} is
\[
  \varrho_{E,\ell}\colon G_\Q \to \GL_2(\Zl).
\]
Its rational version $V_\ell(E) := \Tl(E) \otimes_{\Zl} \Ql$ yields
$\varrho_{E,\ell}\colon G_\Q \to \GL_2(\Ql)$.
\end{definition}

\subsection{Properties of $\varrho_{E,\ell}$}

\begin{theorem}[Weil, 1940]
\label{thm:weil-hasse}
Let $E/\Q$ have good reduction at $p \neq \ell$. Then $\varrho_{E,\ell}$ is
unramified at $p$, and
\[
  \det(1 - \varrho_{E,\ell}(\Frob_p)\,X) = 1 - a_p X + p X^2,
\]
where $a_p = p + 1 - \#E(\Fp)$. In particular,
$|\alpha_p|, |\beta_p| = \sqrt{p}$ where $\alpha_p, \beta_p$ are the roots
(Hasse's theorem: $|a_p| \leq 2\sqrt{p}$).
\end{theorem}

\begin{proof}[Sketch]
The characteristic polynomial of $\Frob_p$ on $V_\ell(E)$ computes the zeta
function of $E$ over $\Fp$ via the Lefschetz trace formula on the Tate module.
Hasse's inequality $|a_p| \leq 2\sqrt{p}$ follows from the positivity of the
Rosati involution.
\end{proof}

\subsection{Ramification at primes of bad reduction}

At a prime $p$ of additive reduction, $\varrho_{E,\ell}|_{I_p}$ is non-trivial and
captures the local Néron model structure. For semi-stable reduction (multiplicative),
the inertia acts unipotently (Tate curve theory).

\subsection{The Weil pairing}

The natural alternating pairing $e_\ell\colon \Tl(E) \times \Tl(E) \to \Tl(\mu)$
gives
\[
  \det(\varrho_{E,\ell}) = \chi_\ell,
\]
the cyclotomic character. This is the key constraint linking the two-dimensional
representation to the one-dimensional cyclotomic character.

%=========================================================
\section{The Modularity Theorem}
\label{sec:modularity}
%=========================================================

\begin{theorem}[Modularity Theorem; Wiles 1995, Taylor--Wiles 1995,
  Breuil--Conrad--Diamond--Taylor 2001]
\label{thm:modularity}
Every elliptic curve $E$ over $\Q$ is modular: there exists a newform
$f = \sum_{n=1}^\infty a_n(f) q^n$ of weight 2 and some level $N$ such that
\[
  a_p(E) = a_p(f) \quad \text{for all primes } p \nmid N.
\]
Equivalently, $L(E, s) = L(f, s)$.
\end{theorem}

\begin{remark}
Wiles proved the semi-stable case (1995), which sufficed to prove Fermat's Last
Theorem. The full theorem for all elliptic curves over $\Q$ was established by
Breuil, Conrad, Diamond, and Taylor (2001).
\end{remark}

\subsection{Strategy of proof}

The proof proceeds via the following steps:

\begin{enumerate}[label=(\roman*)]
\item \textbf{Deformation theory}: For a prime $p | N$ (taken as $p = 3$ or $p = 5$
by Wiles), consider the mod-$p$ representation $\bar{\varrho}_{E,p}$.

\item \textbf{Modularity lifting}: Show that $\varrho_{E,p}$ (the $p$-adic
representation) is modular by lifting the mod-$p$ residual representation
$\bar{\varrho}_{E,p}$ using the \emph{numerical criterion} (the Taylor--Wiles method).

\item \textbf{Taylor--Wiles method}: Construct an auxiliary system of Galois
representations (the \emph{Hecke algebras}) and prove an $R = T$ theorem
(deformation ring $R$ equals Hecke algebra $T$) by a comparison of dimensions
via the Wiles numerical criterion.

\item \textbf{Base case}: Use the known modularity of $E$ modulo small primes via
Langlands--Tunnell (for $p = 3$) or a different base case.
\end{enumerate}

\subsection{Fermat's Last Theorem as corollary}

\begin{corollary}[Fermat's Last Theorem; Wiles 1995]
For $n \geq 3$, the equation $x^n + y^n = z^n$ has no solutions in non-zero integers.
\end{corollary}

\begin{proof}[Proof (given Theorem~\ref{thm:modularity})]
Suppose $a^p + b^p = c^p$ with $abc \neq 0$ and prime $p \geq 5$ (small cases are classical).
Frey (1986) constructed the elliptic curve $E_{a,b,c}\colon y^2 = x(x-a^p)(x+b^p)$.
Ribet (1990) proved (using Mazur's principle) that if $E_{a,b,c}$ were modular,
its associated newform would have level $N = 2$ --- but there are no weight-2 newforms
of level 2 (the relevant space is 0-dimensional). This contradicts modularity, so no
Frey curve exists, hence no Fermat solution.
\end{proof}

%=========================================================
\section{Modular Galois Representations}
\label{sec:modular-forms}
%=========================================================

Let $f = \sum_{n \geq 1} a_n q^n$ be a normalised Hecke eigenform of weight $k$,
level $N$, and Nebentypus $\chi$. Let $K_f = \Q(\{a_n\})$ be its coefficient field.

\begin{theorem}[Deligne 1971]
\label{thm:deligne}
For each prime $\ell$ and embedding $\iota\colon K_f \hookrightarrow \overline{\Q}_\ell$,
there exists a unique (up to isomorphism) irreducible, continuous representation
\[
  \varrho_{f,\ell}\colon G_\Q \to \GL_2(\mathcal{O}_{K_f \otimes \Ql})
\]
unramified outside $N\ell$, such that for all primes $p \nmid N\ell$:
\[
  \tr(\varrho_{f,\ell}(\Frob_p)) = a_p(f), \qquad
  \det(\varrho_{f,\ell}(\Frob_p)) = \chi(p)\, p^{k-1}.
\]
\end{theorem}

\subsection{The Ramanujan--Deligne bound}

\begin{theorem}[Deligne 1974; Weil Conjectures for varieties]
\label{thm:ramanujan-deligne}
For $f$ a weight-$k$ newform as above and $p \nmid N$, the eigenvalues
$\alpha_p, \beta_p$ of $\varrho_{f,\ell}(\Frob_p)$ satisfy
\[
  |\iota(\alpha_p)| = |\iota(\beta_p)| = p^{(k-1)/2}
\]
for every embedding $\iota\colon K_f \hookrightarrow \mathbb{C}$.
For $k = 2$, this gives $|a_p| \leq 2\sqrt{p}$, which is the \emph{Hasse bound}
for elliptic curves (weight-2 forms correspond to elliptic curves via
modularity).  For $k = 12$ and $f = \Delta$, this is Ramanujan's original
conjecture; for general $k$ this is the \emph{Ramanujan--Petersson conjecture}.
\end{theorem}

\begin{proof}[Proof idea]
Deligne realises $\varrho_{f,\ell}$ in the étale cohomology $H^{k-1}_{\text{ét}}$
of a Kuga--Sato variety and applies his proof of the Weil conjectures (Riemann
hypothesis for varieties over finite fields), which gives the bound $p^{(k-1)/2}$
on the eigenvalues of Frobenius.
\end{proof}

%=========================================================
\section{Serre's Modularity Theorem}
\label{sec:serre}
%=========================================================

\begin{definition}
A \textbf{mod-$\ell$ Galois representation} is a continuous homomorphism
$\bar{\varrho}\colon G_\Q \to \GL_2(\overline{\Fp})$.
It is \textbf{odd} if $\det(\bar{\varrho}(\text{complex conjugation})) = -1$,
and \textbf{irreducible} if it has no $G_\Q$-stable line.
\end{definition}

\begin{conjecture}[Serre 1987]
Every odd, irreducible continuous representation
$\bar{\varrho}\colon G_\Q \to \GL_2(\overline{\mathbb{F}}_\ell)$
arises from a newform of precisely the weight $k(\bar{\varrho})$ and level
$N(\bar{\varrho})$ prescribed by Serre's recipe.
\end{conjecture}

\begin{theorem}[Khare--Wintenberger 2009; Serre's Modularity Theorem]
\label{thm:serre}
Serre's conjecture is true: every odd irreducible
$\bar{\varrho}\colon G_\Q \to \GL_2(\overline{\mathbb{F}}_\ell)$
is modular.
\end{theorem}

\begin{proof}[Proof strategy]
Khare and Wintenberger prove this by a two-step induction:
\begin{enumerate}[label=(\alph*)]
\item \textbf{Level-lowering}: Use Ribet's method and modularity lifting theorems
  (in the spirit of Wiles) to reduce to the case of small level and weight.
\item \textbf{Weight-lowering}: Reduce to characteristic-$p$ representations of
  small conductor using $p$-adic Hodge theory and Fontaine--Mazur-type arguments.
\item \textbf{Base case}: Verify the conjecture for small weight and level by an
  explicit finite computation.
\end{enumerate}
The proof uses the full strength of the Taylor--Wiles--Kisin modularity lifting
machinery.
\end{proof}

\begin{corollary}
Serre's Modularity Theorem implies the Modularity Theorem (Theorem~\ref{thm:modularity})
and strengthens the understanding of how residual representations lift to
characteristic-zero representations.
\end{corollary}

%=========================================================
\section{\texorpdfstring{$p$}{p}-adic Representations and Fontaine Theory}
\label{sec:fontaine}
%=========================================================

The local theory of $p$-adic representations (i.e., continuous $G_{\Q_p}$-representations
on $\Q_p$-vector spaces) is controlled by Fontaine's \emph{$p$-adic Hodge theory}.

\subsection{Period rings}

Fontaine constructed several period rings to classify $p$-adic representations:

\begin{definition}
\begin{itemize}
  \item $\BdR$: the de Rham period ring; a complete discrete valuation field with
    residue field $\mathbb{C}_p$ and a filtration.
  \item $\Bcrys \subset \BdR$: the crystalline period ring with Frobenius $\phi$
    and a filtration.
  \item $B_{\mathrm{st}} \supset \Bcrys$: the semi-stable period ring, with an
    additional monodromy operator $N$.
\end{itemize}
\end{definition}

\begin{definition}
A $p$-adic representation $V$ of $G_{\Q_p}$ is:
\begin{itemize}
  \item \textbf{Hodge--Tate} if $V \otimes_{\Q_p} \mathbb{C}_p \cong \bigoplus_i \mathbb{C}_p(i)^{h_i}$,
  \item \textbf{de Rham} if $D_{\mathrm{dR}}(V) := (V \otimes_{\Q_p} \BdR)^{G_{\Q_p}}$ has dimension $n = \dim V$,
  \item \textbf{crystalline} if $D_{\mathrm{crys}}(V) := (V \otimes_{\Q_p} \Bcrys)^{G_{\Q_p}}$ has dimension $n$,
  \item \textbf{semi-stable} if $D_{\mathrm{st}}(V) := (V \otimes_{\Q_p} B_{\mathrm{st}})^{G_{\Q_p}}$ has dimension $n$.
\end{itemize}
\end{definition}

\begin{theorem}[Fontaine--Faltings--Tsuji]
\label{thm:fontaine-comparison}
The $\ell$-adic representation $V_\ell(E)$ of an elliptic curve $E/\Q_p$ is:
\begin{itemize}
  \item Crystalline if $E$ has good reduction at $p$,
  \item Semi-stable if $E$ has semi-stable (multiplicative) reduction at $p$.
\end{itemize}
\end{theorem}

\subsection{Wach modules and integral structures}

For crystalline representations, Wach modules provide integral ($\Zl$-lattice)
versions of the Fontaine classification, allowing the study of representations
over $\Zl$ rather than $\Ql$.

\subsection{The Fontaine--Mazur conjecture}

\begin{conjecture}[Fontaine--Mazur 1993]
An irreducible $p$-adic representation $\varrho\colon G_\Q \to \GL_n(\Ql)$ comes
from geometry (i.e., from the étale cohomology of a variety over $\Q$) if and only
if it is \emph{de Rham} at $p$ (locally at $p$) and unramified at almost all primes.
\end{conjecture}

This conjecture is proved in important special cases (e.g., $n = 2$, connected to
modularity lifting).

%=========================================================
\section{The Langlands Correspondence}
\label{sec:langlands}
%=========================================================

\subsection{The global Langlands correspondence for $\GL_n$}

\begin{conjecture}[Global Langlands correspondence for $\Q$]
\label{conj:langlands}
There is a bijection between:
\begin{enumerate}[label=(\roman*)]
\item $n$-dimensional irreducible continuous representations
  $\varrho\colon G_\Q \to \GL_n(\overline{\Q}_\ell)$ (of geometric origin),
\item Cuspidal automorphic representations $\pi$ of $\GL_n(\mathbb{A}_\Q)$
\end{enumerate}
such that $L(s, \varrho) = L(s, \pi)$ for the associated $L$-functions.
\end{conjecture}

This conjecture remains \textbf{open} for $\Q$ in general. However, special cases
are established:

\begin{itemize}
\item $n = 1$: Class field theory / Artin reciprocity (proved, 20th century).
\item $n = 2$: Weight-$k$ modular forms (proved for elliptic curves by Wiles et al.,
  and more generally for modular Galois representations by Deligne).
\item Function fields: Proved by Lafforgue (2002), see Section~\ref{sec:recent}.
\end{itemize}

\subsection{Local Langlands correspondence}

\begin{theorem}[Local Langlands; Langlands 1970, Henniart 2000, Harris--Taylor 2001]
\label{thm:local-langlands}
For any prime $p$, there is a natural bijection
\[
  \left\{
    \begin{array}{c}
      n\text{-dim. Frobenius-semi-}\\ \text{simple } W_{D}(\Q_p)\text{-reps}
    \end{array}
  \right\}
  \xrightarrow{\;\sim\;}
  \left\{
    \begin{array}{c}
      \text{irreducible smooth reps}\\ \text{of } \GL_n(\Q_p)
    \end{array}
  \right\}
\]
where $W_D(\Q_p)$ denotes the Weil--Deligne group.
\end{theorem}

\subsection{Functoriality}

Langlands's functoriality principle states that a homomorphism of $L$-groups
${}^L G \to {}^L H$ should correspond to a transfer of automorphic representations
from $G(\mathbb{A})$ to $H(\mathbb{A})$. Proved in many cases (base change, symmetric
power lifts for $\GL_2$) but open in general.

%=========================================================
\section{Recent Results}
\label{sec:recent}
%=========================================================

\subsection{Lafforgue's theorem over function fields}

\begin{theorem}[Lafforgue 2002; Fields Medal]
\label{thm:lafforgue}
The global Langlands correspondence holds for $\GL_n$ over a function field
$\mathbb{F}_q(C)$ (the function field of a curve $C$ over a finite field $\mathbb{F}_q$):
there is a bijection between $n$-dimensional $\ell$-adic representations of the
absolute Galois group of $\mathbb{F}_q(C)$ and cuspidal automorphic representations
of $\GL_n(\mathbb{A}_{\mathbb{F}_q(C)})$, matching $L$-functions.
\end{theorem}

This was the first complete proof of the global Langlands correspondence for $\GL_n$
with $n > 2$. It uses Drinfeld's work for $\GL_2$ (1974, 1980) as foundation.

\subsection{Potential automorphy}

\begin{theorem}[Taylor 2004; Harris--Shepherd-Barron--Taylor 2010]
\label{thm:potential-automorphy}
Let $\varrho\colon G_\Q \to \GL_n(\overline{\Q}_\ell)$ be a "regular, odd"
$\ell$-adic representation. Then $\varrho$ is \emph{potentially automorphic}:
there exists a totally real field $F$ such that $\varrho|_{G_F}$ is automorphic.
\end{theorem}

Potential automorphy suffices to prove that $L(\varrho, s)$ has analytic continuation
and satisfies a functional equation, even when full automorphy is not established.

\subsection{The Taylor--Wiles--Kisin method}

Diamond (1997), Kisin (2009), and many others extended the Taylor--Wiles method to
cover all 2-dimensional de Rham representations of $G_{\Q_p}$. Kisin's proof uses
\emph{Breuil--Kisin modules} (analogues of Wach modules) and resolves most cases of
the Fontaine--Mazur conjecture for $n = 2$.

\subsection{$p$-adic Langlands correspondence}

For $\GL_2(\Q_p)$, Colmez (2010) established a $p$-adic local Langlands correspondence:
a bijection between 2-dimensional trianguline $p$-adic representations of $G_{\Q_p}$
and locally analytic representations of $\GL_2(\Q_p)$, compatible with the classical
local Langlands correspondence.

\subsection{Higher-dimensional and geometric advances}

Recent work by Scholze (using perfectoid spaces and diamonds), V.~Lafforgue (2018, for
reductive groups over function fields), and Fargues--Scholze (2021, geometrisation of
local Langlands) has dramatically advanced the structural understanding. Fargues and
Scholze have geometrised the local Langlands programme, constructing a sheaf-theoretic
version over the Fargues--Fontaine curve.

%=========================================================
\section{References}
%=========================================================

\begin{thebibliography}{99}

\bibitem{Wiles1995}
A.~Wiles,
\textit{Modular elliptic curves and Fermat's last theorem},
Annals of Mathematics \textbf{141} (1995), 443--551.

\bibitem{TaylorWiles1995}
R.~Taylor and A.~Wiles,
\textit{Ring-theoretic properties of certain Hecke algebras},
Annals of Mathematics \textbf{141} (1995), 553--572.

\bibitem{BCDT2001}
C.~Breuil, B.~Conrad, F.~Diamond, and R.~Taylor,
\textit{On the modularity of elliptic curves over $\mathbb{Q}$: Wild 3-adic exercises},
Journal of the American Mathematical Society \textbf{14} (2001), 843--939.

\bibitem{KhareWintenberger2009}
C.~Khare and J.-P.~Wintenberger,
\textit{Serre's modularity conjecture (I)},
Inventiones Mathematicae \textbf{178} (2009), 485--504.

\bibitem{Serre1987}
J.-P.~Serre,
\textit{Sur les représentations modulaires de degré 2 de $\mathrm{Gal}(\overline{\mathbb{Q}}/\mathbb{Q})$},
Duke Mathematical Journal \textbf{54} (1987), 179--230.

\bibitem{Fontaine1994}
J.-M.~Fontaine,
\textit{Représentations $p$-adiques semi-stables},
Astérisque \textbf{223} (1994), 113--184.

\bibitem{Lafforgue2002}
L.~Lafforgue,
\textit{Chtoucas de Drinfeld et correspondance de Langlands},
Inventiones Mathematicae \textbf{147} (2002), 1--241.

\bibitem{Deligne1974}
P.~Deligne,
\textit{La conjecture de Weil. I},
Publications Mathématiques de l'IHÉS \textbf{43} (1974), 273--307.

\bibitem{Langlands1970}
R.~P.~Langlands,
\textit{Problems in the theory of automorphic forms},
Lectures in Modern Analysis and Applications III,
Lecture Notes in Mathematics \textbf{170}, Springer, 1970, 18--61.

\bibitem{HarrisTaylor2001}
M.~Harris and R.~Taylor,
\textit{The geometry and cohomology of some simple Shimura varieties},
Annals of Mathematics Studies \textbf{151}, Princeton University Press, 2001.

\end{thebibliography}

\end{document}
